\documentclass[11pt]{amsart}
\usepackage{geometry}
\geometry{letterpaper}
\usepackage{graphicx}
\usepackage{amssymb}
\usepackage{epstopdf}
\DeclareGraphicsRule{.tif}{png}{.png}{`convert #1 `dirname #1`/`basename #1 .tif`.png}

\title{ Problema 2.10}
\author{}
%\date{}

\begin{document}
\maketitle

Dado como datos tres números reales, identifique cuál es el mayor. Considere que los números pueden ser iguales. Desarrolle el diagrama de flujo correspondiente.\\

\textbf{DATOS:} $A$, $B$, $C$ (variables de tipo real).\\

\begin{center}
\begin{tabular}{|c|c|c|c|}
\hline
\multicolumn{3}{ |c| }{Tabla 2.15} \\ \hline
  \hline
  NUMERO DE CORRIDAD & DATO CAL & RESLTADO \\
  \hline
  1 & 8.75 & "Aprobado"  \\
  \hline
  2 & 7.90 &  "Reprobado"  \\
  \hline
  3 & 8.00 &  "Aprobado" \\
  \hline
  4 & 9.50 & "aprobado" \\
  \hline
  5 & 8.35 & "Aprobado" \\
  \hline
\end{tabular}
\end{center}

\end{document}
